\expandafter\def\csname HistoricalNoteText536\endcsname{Seit kurzem ist dieser mit einem Stahldraht vertauscht.}
\SourceLeaf{532}in Leipzig und Berlin und beinahe vollständig in Braunschweig angestellt; ausserdem auch zur Hälfte in Copenhagen, wo durch Versehen der 24. und 25. September anstatt des 23. und 24. genommen wurden. Die vollständige Bekanntmachung dieser Beobachtungen würde jedoch geringeres Interesse haben, da der Verlauf an diesen beiden Tagen sehr regelmässig war, obgleich mehrere an sich sehr kleine Anomalien in den ersten 24 Stunden an allen vier Plätzen eine bewundrungswürdige Harmonie gezeigt haben. Merkwürdig bleibt indessen, dass, einer Zeitungsnachricht zufolge, am 23. September Abends in Glasgow ein sehr starkes Nordlicht gesehen worden ist, welches mithin ganz entschieden, wenigstens keinen sich bis Norddeutschland erstreckenden Einfluss auf die Magnetnadel gehabt hat.

Die Anwesenheit des Herrn Prof.\ \textsc{Weber} in Leipzig veranlasste inzwischen, noch einige ausserordentliche Stunden zu gleichzeitigen Beobachtungen in Göttingen und Leipzig festzusetzen, wozu die Tage 1. und 2. October Morgens $7\tfrac12$ bis $8\tfrac12$, Mittags $12\tfrac12$ bis $1\tfrac12$, und Abends 8 bis 10 Uhr gewählt wurden. Abgesehen von einigen kleinen Versäumnissen wurden diese Stunden an beiden Orten inne gehalten; im hiesigen magnetischen Observatorium beobachtete mein Sohn, \textsc{Wilhelm Gauss}, in Leipzig Herr Prof.\ \textsc{Weber}, Herr Prof.\ \textsc{Möbius} und Herr Dr.\ \textsc{Thieme}. — — — Man wird nicht ohne Vergnügen die grosse Übereinstimmung nicht blos in den grossen Bewegungen, welche am Abend des 1. October stattfanden, sondern fast in sämmtlichen kleinen bemerken, so dass deren Quellen sich als auf grosse Ferne hinwirkende, obwohl zur Zeit noch sehr räthselhafte Kräfte, auf das Unverkennbarste ausweisen. In Leipzig waren die Anomalien im Allgemeinen etwas kleiner als in Göttingen; letzterem Orte wird daher der Heerd der wirkenden Kräfte näher gewesen sein. Ich bemerke nur noch, dass während eines Theils jener Stunden ich selbst an einem zweiten in der hiesigen Sternwarte aufgestellten Apparat, wovon ich bald eine ausführlichere Nachricht zu geben gedenke, beobachtet habe, und dass diese Beobachtungen einen fast vollkommenen Parallelismus mit denen des hiesigen magnetischen Observatoriums in den grösseren und kleineren Bewegungen ergeben haben; ein ähnlicher Erfolg hatte auch am 23. und 24. September, so wie bei vielen sonstigen Versuchen, statt, in dem Maasse, dass schon öfters die Uhren an beiden Plätzen blos mittelst der magnetischen Erscheinungen bis auf einen kleinen Bruchtheil einer Zeitminute genau verglichen werden konnten. Dasselbe gelingt mittelst
\SourceLeaf{533}der grösseren Bewegungen am 1. und 2. October zwischen Göttingen und Leipzig, wo an beiden Orten, die Uhren nur wenige Secunden von der mittleren Ortszeit abwichen.

Durch diese Erfahrungen erhalten nun auch die kleinen, in sehr kurzen Zeitfristen wechselnden Schwankungen der Magnetnadel ein überaus grosses Interesse; man muss wünschen, dass auch diese durch die Beobachtungen an vielen von einander entfernten Plätzen sorgfältig verfolgt werden, und es wird daher unumgänglich nöthig alle Beobachtungen in recht kurzen Zeitintervallen zu machen. Bisher beobachteten wir von 5 zu 5 Minuten: aber auch dieses Intervall ist noch fast zu lang, und wir denken künftig immer von 3 zu 3 Minuten den Stand der Magnetnadel an den verabredeten Tagen zu bestimmen. Ich darf dabei nicht unbemerkt lassen, dass das Verfahren, welches der Herr Herausgeber dieser Annalen (Bd.\ XXXII. S.\ 569 bis 572) erklärt hat, von uns nur anfänglich gebraucht, aber schon lange mit einem etwas abgeänderten vertauscht ist. Um den Stand der Magnetnadel für einen Augenblick zu erhalten, beobachten wir sie in sechs verschiedenen, immer um Eine Schwingungsdauer getrennten Momenten, und so, dass der gewünschte Moment in die Mitte fällt. Anstatt der genauen Schwingungsdauer wird die nächste runde Zahl von Secunden (oder vielmehr von Uhrschlägen) gewählt, z.\ B. im magnetischen Observatorium $20''$ anstatt $20''4$. Die Beobachtungen am 2. October für $8^{\mathrm h}\,15'$ Abends standen daher so:
\begin{center}
\begin{tabular}{rrr|r}
$8^{\mathrm h}$ & $14'$ & $10''$ & $672.6$\\
 & & $30$ & $672.3$\\
 & & $50$ & $671.3$\\
 & $15$ & $10$ & $671.8$\\
 & & $30$ & $669.9$\\
 & & $50$ & $670.8$
\end{tabular}
\end{center}
\Needspace{13\baselineskip}
Hieraus ergeben sich fünf Mittel, die eigentlich den beigesetzten Zeiten correspondiren:
\begin{center}
\begin{tabular}{rrr|r}
$8^{\mathrm h}$ & $14'$ & $20''$ & $672.45$\\
 & & $40$ & $671.80$\\
 & $15$ & $0$ & $671.55$\\
 & & $20$ & $670.85$\\
 & & $40$ & $670.35$
\end{tabular}\\[2mm]
\begin{tabular}{lrr}
und daraus das Mittel für & $8^{\mathrm h}\,15'$ & $671.40$
\end{tabular}
\end{center}

\SourceLeaf{534}Ich habe absichtlich dieses Beispiel gewählt, wo die Nadel schnelle Veränderungen zeigte. die selbst von 20 zu 20 Secunden sich so entschieden darstellen. Wir haben Fälle genug, wo ein ähnlicher Erfolg selbst in halb so grossen Zeitintervallen eintritt. Gewisse Abänderungen in jener Beobachtungsart (die wir öfters anwenden) zu erklären, so wie die Rechtfertigung jener Art das Mittel zu nehmen, die mit gutem Vorbedacht gewählt ist, muss ich mir für eine andere Gelegenheit vorbehalten. Aber unerwähnt lassen darf ich nicht (da der Herr Herausgeber dieser Annalen a.\ a.\ O. es nicht ausdrücklich bemerkt hat), dass es eine wesentliche Bedingung für die Zulässigkeit aller dieser Beobachtungsarten ist, die Nadel vorher so viel wie möglich beruhigt zu haben, so dass die Schwingungen nur eine geringe Anzahl von Scalentheilen betragen. Im hiesigen magnetischen Observatorium ist eine solche Beruhigung oder wenigstens eine Wiederholung derselben, im Laufe der Beobachtung selten nöthig. Wer aber in einem weniger günstigen Local beobachtet, darf durchaus nicht unterlassen, dies, so oft es nöthig wird, in der Zwischenzeit mit den bekannten Mitteln zu thun.

Da bei der gegenwärtig als nothwendig sich zeigenden Verengerung der Zwischenzeiten die Beobachtungen sehr viel mühsamer werden als früher, wo die Forderung sich auf die Aufzeichnung von Stunde zu Stunde beschränkte, so ist mehrseitig der Wunsch geäussert, künftig sowohl die Anzahl als die Dauer der Termine etwas zu verkürzen. — — —

Göttingen, den 5. November 1834.

\begin{samepage}
\bigskip
\pdfbookmark[0]{Magnetische Beobachtungen. 1835}{magnetische_beobachtungen_1835}
\begin{center}
\hrule\vspace{2mm}
Göttingische gelehrte Anzeigen. 1835 März 7.
\vspace{2mm}\hrule\vspace{5mm}
\end{center}
In der Sitzung der Königl.\ Societät am 14. Februar stattete der Hofr.\ \textsc{Gauss} einen Bericht über die in dem magnetischen Observatorium, und in Verbindung damit anderwärts gemachten Beobachtungen ab, woraus wir hier einen Auszug mittheilen, der als
\begin{center}\itshape eine Fortsetzung der am 9. August 1834 gegebenen Nachricht\end{center}
\noindent betrachtet werden kann.
\end{samepage}

\Needspace{13\baselineskip}
\SourceLeaf{535}Die täglich zweimaligen Aufzeichnungen des Standes der Nadel sind ununterbrochen fortgesetzt, und umfassen nun bereits beinahe ein volles Jahr. Die monatlichen Mittel, seit Julius v.\ J., waren:
\begin{center}
\begin{tabular}{rl|rrr|rrr}
 & & \multicolumn{3}{c|}{8 Uhr Vorm.} & \multicolumn{3}{c}{1 Uhr Nachm.}\\\hline
1834 & August & $18^0$ & $38'$ & $48''1$ & $18^0$ & $49'$ & $11''0$\\
 & September & & $36$ & $58.4$ & & $46$ & $32.3$\\
 & October & & $37$ & $18.4$ & & $44$ & $47.2$\\
 & November & & $37$ & $38.4$ & & $43$ & $4.3$\\
 & December & & $37$ & $54.8$ & & $41$ & $32.7$\\
1835 & Januar & & $37$ & $51.5$ & & $42$ & $14.4$
\end{tabular}
\end{center}

Die verabredeten Beobachtungen an bestimmten Tagen in kurzen ununterbrochenen Zeitfristen, mit deren Einrichtung in den letzten Monaten einige an einem andern Orte bekannt gemachte Abänderungen getroffen sind, haben seit der letzten Nachricht an vier Hauptterminen Statt gefunden, einige ausserordentliche Nebentermine ungerechnet. Die Theilnahme an denselben hat sich bereits weiter ausgebreitet, und wird bald noch weiter verbreitet werden, auch sind daraus schon sehr merkwürdige Resultate hervorgegangen, denen ähnlich, welche in dem frühern Bericht erwähnt wurden. Eine graphische Darstellung der Harmonie unter den Beobachtungen vom 1. und 2. October, und vom 29. und 30. November in Göttingen, Leipzig und Berlin, wird nächstens in \textsc{Poggendorffs} Annalen der Physik erscheinen: noch merkwürdiger aber ist die Übereinstimmung der Beobachtungen vom 5. und 6. November in Copenhagen und Mailand in allen zahlreichen und auffallend grossen Schwankungen, von welchen gleichfalls eine Zeichnung an einem andern Orte gegeben werden wird. Wir treten hier in eine Welt von geheimnissvollen Naturkräften, deren wunderbar wechselndes Spiel sich über den halben Durchschnitt von Europa, in gleichem Augenblick, und bis in die kleinsten Nuancen auf gleiche Weise, offenbart, und deren Wirkungskreis zu ermessen diese Standlinie noch viel zu klein erscheint.

Die hiesigen Einrichtungen für magnetische Beobachtungen haben inzwischen mehrere wesentliche Erweiterungen erhalten. Für manche Beobachtungen ist, wenn grosse Schärfe verlangt wird, die Zuziehung eines zweiten Apparats, in einiger Entfernung vom Hauptapparate, unumgänglich nothwendig, um von den stündlichen Veränderungen der magnetischen Kraft Rechnung tragen zu können, Zu diesem Zweck ist seit August v.\ J., nachdem die im Jahre 1832 ge\SourceLeaf{536}brauchten Apparate an das physicalische Cabinet abgegeben sind, in der Sternwarte ein grosser Magnetstab aufgehängt, mit übrigens ganz ähnlichem Zubehör, wie der Stab im magnetischen Observatorium. Der Magnetstab in der Sternwarte, gleichfalls aus Uslarschem Gussstahl, ist 4 Fuss lang, fast drei Zoll breit und über einen halben Zoll dick, und wiegt 25 Pfund. Er hängt an einem 16 Fuss langen tausendfachen Seidenfaden\HistoricalNote{536}, der oberhalb der Decke des Saals seine Befestigung hat, und durch eine kleine in dieser Decke gemachte Öffnung frei durchgeht. Der nächste Grund zur Wahl eines so schweren Stabes war die Absicht, den Luftzug, welcher in diesem Local nicht immer ganz abgehalten werden kann, und der auf die kleinern Apparate, ungeachtet der Beschützung durch einen umschliessenden Kasten öfters störend einwirkte, unschädlich zu machen. Der Erfolg hat nicht nur \textit{dieser} Erwartung entsprochen, sondern auch die andern rücksichtlich der Genauigkeit aller daran zu machenden Beobachtungen noch weit übertroffen. Nur absolute Beobachtungen der Declination und Intensität bleiben natürlich wegen des in der Sternwarte vielfach vorhandenen Eisens davon ausgeschlossen.

Die grösste Schwingung, welche der den Stab einschliessende Kasten verstattet, beträgt etwa 27 Grad; die grösste, welche auf der Scale unmittelbar noch gemessen werden kann, 9 bis 10 Grad, indem bei grössern die Gesichtslinie des Fernrohrs nicht mehr auf den fast vier Zoll breiten Spiegel trifft. Ist der Stab einmal in Schwingungen gesetzt, so nehmen diese in geometrischer Progression so langsam ab, dass sie oft erst nach 10 oder mehreren Stunden auf die Hälfte herabkommen, obwohl zuweilen auch viel früher, von welchem Umstande unten noch besonders die Rede sein wird. Die Dauer einer Schwingung des jetzt eingehängten Stabes, des stärksten aus einer grössern Zahl, die für das physicalische Cabinet angefertigt sind, beträgt etwas über 42 Secunden, und diese Grösse, welche wegen Temperatur und Veränderlichkeit des Erdmagnetismus einigen, obwohl sehr kleinen Veränderungen unterworfen ist (so wie auch vielleicht im Laufe der Zeit eine bis jetzt noch gar nicht spürbare Veränderung der Kraft des Stabes selbst eintreten kann), wird aus einigen wenigen Schwingungen schon so scharf bestimmt, dass man dann den Stab auf 8 und mehrere Stunden verlassen kann, ohne nachher über die Anzahl der inzwischen vollendeten Schwingungen zweifelhaft zu bleiben.

\SourceLeaf{537}Eben so interessant, wie die rein magnetischen Beobachtungen sind die mit diesem Apparat anzustellenden electrodynamischen Versuche. Zu diesem Zweck ist der Stab von einem ähnlichen Multiplicator umgeben, wie der Stab des magnetischen Observatoriums, nur dass jener grössere Dimensionen, und eine Drahtlänge von 2700 Fuss in 270 Umwindungen hat. Dieser Multiplicator ist in die grosse schon in dem frühern Bericht erwähnte Drahtkette gebracht, welche die Sternwarte, das magnetische Observatorium und das physicalische Cabinet verbindet, und in welcher der galvanische Strom zusammen eine Drahtlänge von 11000 Fuss, also fast einer halben geographischen Meile zu durchlaufen hat, und dann drei magnetische Apparate zugleich afficirt, nemlich
\begin{list}{}{\setlength{\leftmargin}{2em}\setlength{\labelwidth}{2em}\setlength{\itemsep}{0pt}\setlength{\parsep}{0pt}}
\item[I.] den 25pfündigen Stab in der Sternwarte,
\item[II.] den 4pfündigen Stab im magnetischen Observatorium\newline (Multiplicator von 200 Umwindungen)
\item[III.] den einpfündigen Stab im physikalischen Cabinet\newline (Multiplicator von 160 Umwindungen).
\end{list}
\noindent Einzelne Theile der Kette können in vielfachen Combinationen nach Gefallen mit Leichtigkeit abgesperrt werden.

Von den zahlreichen Versuchen, welche schon jetzt mit diesen Apparaten gemacht sind, führen wir hier nur einige an.

Wenn ein galvanischer Strom mit der Kette in Verbindung gesetzt wird, so erscheinen die Bewegungen der Magnetstäbe in den drei Apparaten so augenblicklich, dass ihr Anfang sich auf einen kleinen Bruch einer Zeitsecunde genau beobachten lässt. Die Vergleichung der Uhren bei den drei Apparaten liefert so vollkommen übereinstimmende Resultate, der Strom möge an dem einen Ende, oder an dem andern, oder in der Mitte erzeugt sein, dass daraus die Unmessbarkeit der Zeit, in welcher der Strom eine halbe Meile durchläuft, vollkommen bestätigt wird. Nach den interessanten Versuchen von \textsc{Wheatstone}, welche neuerlich in den \textit{Philosophical Transactions} für 1834 bekannt gemacht sind, und nach welchen der electrische Strom im Metall eine grössere Geschwindigkeit zu haben scheint, als das Licht im Raume, liess sich freilich ein solcher Erfolg schon vermuthen, obwohl sich daraus doch nicht unbedingt auf das Verhalten eines \textit{galvanischen} Stroms, und dessen Einwirkung auf die Magnetnadel schliessen liess.

Die Intensität eines galvanischen Stroms wird durch die Ablenkung der Magnetnadel, also zunächst durch Scalentheile gemessen oder bestimmt, allein
\SourceLeaf{538}offenbar in den drei Apparaten mit verschiedenen Einheiten, welche von den Dimensionen der Multiplicatoren und der Geltung der Scalentheile in Bogensecunden abhangen. Nun zeigen aber zahlreiche angestellte Versuche, dass zwischen den Ablenkungen an den drei Apparaten durch denselben Strom in einerlei Augenblick stets genau ein constantes Verhältniss Statt findet, der Strom möge an dem einen, oder an dem andern Ende, oder in der Mitte erzeugt sein. Es ergibt sich daraus das wichtige Resultat, dass der Strom in seiner ganzen Länge dieselbe Intensität hat, wenigstens nichts merkliches davon verliert. Man wird in Zukunft besonders aufmerksam darauf sein, ob dieses Resultat auch unter eigenthümlichen Umständen, namentlich während starken Regens, seine Gültigkeit behält.

Bei allen drei Apparaten sind Commutatoren (Gyrotrope) mit der Kette verbunden, wodurch man die Richtung des Stroms mit Leichtigkeit umkehren kann. Dem Commutator in der Sternwarte hat der Hofr.\ \textsc{Gauss} eine eigenthümliche Einrichtung gegeben, wonach diese Umkehrung durch einen einzigen Druck mit dem Finger, also ganz augenblicklich, bewirkt wird. Wenn man diese Umkehrung, immer in so grossen Zeitfristen wie die Schwingungsdauer des Einen Stabes, wiederholt ausführt, so werden die Schwingungen dieses Stabes immer grösser. Man hat dies zu einem Experiment benutzt, wobei eine auffallende mechanische Wirkung hervorgebracht wird. Herr Prof.\ \textsc{Weber} liess zur Seite des Magnetstabes im physikalischen Cabinet eine leichte Auslösung für einen Wecker oder eine Pendeluhr anbringen. Dieses Auslösen gelingt jedesmal durch den von der Sternwarte aus geleiteten Strom nach ein Paar Schwingungen auf das vollkommenste. Dass man mit dem 25pfündigen Stabe eine noch viel stärkere mechanische Wirkung würde hervorbringen können, leuchtet von selbst ein.

Besonders wichtige Dienste leisten diese Apparate bei der Erforschung der mathematischen Gesetze, nach welchen sich die Erzeugung und die Wirkung der von \textsc{Faraday} entdeckten magneto-electrischen Induction richten, und ihrer Zurückführung auf absolute Maasse, worüber der Hofr.\ \textsc{Gauss} den Erfolg seiner Untersuchungen zu seiner Zeit an einem andern Orte bekannt machen wird. Von den dabei angewandten Vorrichtungen erwähnen wir hier nur einer, womit diese Induction auf eine eben so einfache als scharf messbare Art dargestellt wird. Um eine hölzerne Rolle ist ein übersponnener Draht mit 1050 Umwindungen geführt, dessen Enden durch den Commutator mit der Kette in Verbindung gebracht wer\SourceLeaf{539}den. Diese Rolle kann über die freistehende Hälfte eines starken Magnetstabes geführt werden, und während dieser Operation geht allemal durch die Kette ein galvanischer Strom, ein starker, aber von kurzer Dauer, oder ein schwächerer von längerer Dauer, je nachdem die Manipulation schneller oder langsamer geschieht, so dass die Gesammtwirkung Eines Aufschiebens von der Schnelligkeit der Operation unabhängig ist. Der Strom an sich dauert immer nur so lange, wie die Bewegung der Rolle. Das Abziehen der Rolle bringt einen entgegengesetzten Strom hervor, eben so das Aufschieben mit dem entgegengesetzten Ende. Geschieht die Bewegung sehr schnell, so ist die Wirkung des Stroms auf die Magnetnadel in einem der mit der Kette verbundenen Multiplicatoren einem augenblicklichen Stosse von bestimmter Stärke gleich zu setzen. Abziehen und verkehrt wieder Aufstecken bewirkt also zwei gleichnamige Impulse der Magnetnadel, und ein neues Abziehen und wieder umgekehrt Aufschieben würde daher zwei unter sich gleiche, aber den vorigen entgegengesetzte Impulse hervorbringen; allein wenn dazwischen der Commutator gewechselt ist, so geschehen auch die letzten beiden Wirkungen in demselben Sinn, wie die beiden ersten. Ein solcher vollständiger Wechsel (Abziehen, Verkehrtaufstecken und Commutatorumstellung) geschieht ganz bequem in zwei Secunden, und man kann daher, wenn man will, während einer Schwingungsdauer des grossen Magnetstabes bequem und tactmässig 21 Wechsel vollenden, und dadurch letztern in so starke Bewegung bringen, dass die ganze Scale aus dem Gesichtsfelde des Fernrohrs geht. Diese Andeutung wird hinreichen zu übersehen wie die Stärke des durch diese Inductionsart entstehenden galvanischen Stroms mit Schärfe gemessen werden kann. Diese Stärke hängt aber zugleich von dem Widerstande ab, welchen die Kette selbst darbietet, und nimmt mehr oder weniger zu, je nachdem mehr oder weniger Stücke der Kette abgesperrt werden. Auf diese Weise ist das Verhältniss des Widerstandes in den einzelnen Bestandtheilen der Kette und den Multiplicatoren mit grosser Schärfe bestimmt, und durch mannigfaltige Combinationen das schöne von \textsc{Ohm} aufgestellte Gesetz, welches die Intensität eines Stroms bei einer Theilung befolgt, auf das vollkommenste bestätigt. Nahe übereinstimmende Resultate sind auch mit hydrogalvanischen Strömen gefunden; indessen eignen sich diese, wegen der Veränderlichkeit ihrer Stärke weniger zu solchen Bestimmungen, und erfordern jedenfalls deshalb noch besondere Vorsichtsmaassregeln bei den Versuchen. Vielleicht ist nicht uninteressant, wenn hier bemerkt
\SourceLeaf{540}wird, dass der ganze Widerstand in der in der Luft geführten doppelten Drahtverbindung zwischen der Sternwarte und dem physikalischen Cabinet, in einer Drahtlänge von mehr als 6000 Fuss nur ungefähr halb so gross ist, als der Widerstand, welchen der Strom bloss in dem Multiplicator des magnetischen Observatoriums (Drahtlänge 1100 Fuss) findet, oder nur den sechsten Theil des Widerstandes in der ganzen Kette beträgt: indessen erklärt sich dies leicht aus der ungleichen Dicke des Drahts, und alle Versuche bestätigen, dass bei Drähten von einerlei Metall der Widerstand immer im geraden Verhältniss der Länge und im umgekehrten der \textit{Fläche} des Querschnitts steht.

Wir haben oben erwähnt, dass die Abnahme des Schwingungsbogens bei der grossen Nadel in verschiedenen Zeiten sehr ungleich gewesen ist. Ähnliche Verschiedenheiten hatten sich schon im Jahr 1832 bei den kleinen Apparaten gezeigt, auch später bei der Nadel im magnetischen Observatorium: allein diese Verschiedenheiten bleiben immer innerhalb viel engerer Grenzen, als bei dem Stabe der Sternwarte, wo die Abnahme des Schwingungsbogens von einer Schwingung zur folgenden in verschiedenen Versuchsreihen zwischen $\frac{1}{5000}$ und $\frac{1}{50}$ schwankte. Diese merkwürdige Erscheinung hat die Aufmerksamkeit des Hofr.\ \textsc{Gauss} besonders auf sich gezogen, und es scheint dabei ein Zusammentreffen \textit{mehrerer} Ursachen Statt zu finden, die zum Theil noch jetzt räthselhaft bleiben: inzwischen ist es dem Hofr.\ \textsc{Gauss} gelungen, diejenige Ursache, welche bei weitem den stärksten Einfluss hat, auszumitteln. Er bemerkte nemlich, dass allemal der Schwingungsbogen viel schneller abnahm, wenn die Kette geschlossen, als wenn sie offen war, und so war es leicht, als Ursache jener schnellen Abnahme, die Reaction eines in der Kette durch die Schwingung der Nadel selbst, vermöge der Induction, erzeugten galvanischen Stroms zu erkennen, welcher bei der folgenden Rückschwingung die entgegengesetzte Richtung hat, und stets auf Verminderung des Schwingungsbogens wirkt. Diese Erklärung bestätigte sich vollkommen, indem die Abnahme des Schwingungsbogens am langsamsten war bei offner Kette, schneller bei geschlossener aber vollständiger Kette; noch schneller, wenn einzelne Stücke der Kette abgesperrt waren; und am allerschnellsten (so dass der Schwingungsbogen in einer halben Stunde auf die Hälfte kam), wenn die Kette gleich hinter dem Multiplicator des grossen Stabes geschlossen war. Ja diese Unterschiede richteten sich vollkommen nach der Grösse des wirksam bleibenden Theils der Kette.

\SourceLeaf{541}Nachdem diese Erklärung gefunden war, war es leicht, den Erfolg einiger Versuche vorauszusehen, welche wohl zu den auffallendsten im Gebiet des Electromagnetismus gerechnet werden dürfen, und selbst die quantitativen Verhältnisse der Erscheinungen im Voraus zu berechnen, welche auch bei den wiederholt angestellten Versuchen stets auf das vollkommenste bestätigt sind. Es sind folgende.

Wenn der Magnetstab in der Sternwarte (I) in Schwingungen gesetzt wird, etwa so grosse wie der Kasten verstattet, so haben diese gar keinen Einfluss auf die Nadeln im magnetischen Observatorium (II) oder im physikalischen Cabinet (III), sondern diese \textit{bleiben} in Ruhe, wenn sie vorher in Ruhe waren, vorausgesetzt, dass die Kette offen, oder wenigstens die die letzten Nadeln einschliessenden Multiplicatoren davon abgesperrt sind. Allein in dem Augenblick, wo die Kette geschlossen oder z.\ B.\ der Multiplicator von II in die geschlossene Kette hineingebracht wird, fängt die Nadel II sogleich an mitzuschwingen. Ist die Nadel II schon vorher in Schwingung gewesen, so erhalten die Schwingungen den eigenthümlichen Character \textit{gemischter} Schwingungen, wovon die eine von dem Initialzustande abhängt, und dieselbe Periode hat, wie die Schwingungen dieser Nadel unter dem blossen Einfluss des Erdmagnetismus ($20''$), während die andere eine Periode von $42''$ befolgt (wie die grosse Nadel I), und ihre Grösse dem Schwingungsbogen von I proportional ist (etwa $\frac{1}{200}$, wenn die Kette hinter dem Multiplicator von II abgesperrt ist). Dies ist vollkommen mit den Resultaten der Theorie in Übereinstimmung, eben so wie der stets genau bestätigte Umstand, dass die Schwingungen von I und die inducirten Schwingungen von II, obwohl Perioden von gleicher Dauer, doch nicht gleichem Anfang haben, sondern stets eine halbe Schwingungszeit ($21''$) in dieser Beziehung differiren, und zwar in dem Sinn, wie es nach den Statt findenden Umständen die Theorie vorausbestimmt. Was hier beispielsweise von der Nadel II gesagt ist, findet auf ganz ähnliche Weise bei der Nadel III Statt, deren natürliche Schwingungsdauer $14''$ beträgt, und die unter der Einwirkung der Induction zusammengesetzte Schwingungen von $14''$ und $42''$ Periode befolgt.

Ein ganz anderer Erfolg muss der Theorie zufolge in dem Fall Statt finden, wenn eine zweite Nadel, deren natürliche Schwingungsdauer genau eben so gross ist, wie die des grossen Magnetstabes, mit einem Multiplicator sich in der Kette befindet, in welcher der grosse Stab schwingt. Jene, so lange vollkommen ruhig,
